\documentclass[10pt,a4paper]{amsart}
\usepackage[margin=19mm]{geometry}
\usepackage{amsmath,amssymb,mathtools}
\usepackage[scr=rsfs,cal=euler]{mathalfa}
\usepackage{xfrac}
\usepackage[unicode,hidelinks,bookmarks=false]{hyperref}
\newcommand{\ie}{i.e.\ }
\newcommand{\cf}{cf.\ }
\newcommand{\resp}{respectively\ }
\newcommand{\Subsection}{\subsection}
\providecommand{\nobreakdash}{\nobreak\hbox{-}}
\setlength{\emergencystretch}{2em}
\multlinegap0pt
\usepackage{amssymb}
\usepackage{xcolor}
\usepackage{enumerate}
\usepackage{bm}
\usepackage[normalem]{ulem}
\usepackage{tikz}
\newtheorem{theorem}{Theorem}
\DeclareMathOperator{\afs}{AFS}
\newcommand{\aht}{\mathnormal{\mathsf{AHT}}}
\newcommand{\weiseq}{\mathop{\equiv_{\mathrm{s}\mathsf{W}}}}
\newcommand{\weisleq}{\mathop{\leq_{\mathrm{s}\mathsf{W}}}}
\newcommand{\zrt}{\mathnormal{\mathbb Z\text{-}\mathsf{RT}}}
\title{The adjacent Hindman's theorem and the $\mathbb Z$-Ramsey's theorem}
\author{Bruno Fernando Aceves-Mart\'{\i}nez, David J. Fern\'andez-Bret\'on, L. F. Romero-Garc\'{\i}a, Luis F. Villag\'omez-Canela}
\date{Source: arXiv:2412.14558v3 (CC BY 4.0)}
\begin{document}
\maketitle

\noindent\emph{Source and licence.} The adjacent Hindman's theorem and the $\mathbb Z$-Ramsey's theorem. Version arXiv:2412.14558v3 (CC BY 4.0), distributed by arXiv under the Creative Commons Attribution 4.0 International licence, http://creativecommons.org/licenses/by/4.0/, as recorded in the arXiv licence metadata for this exact version and shown on the abstract page https://arxiv.org/abs/2412.14558v3. Full licence notice: \texttt{licenses/arxiv-2412.14558-CC-BY-4.0.txt}.\par\smallskip

\noindent\textit{Editorial scope.} Counted unit: the article's theorem higher-zrtversusaht (source0.tex lines 308--310). The article's complete original proof of it is delivered with it (source0.tex lines 312--339, one \texttt{proof} environment of 28 lines). No mathematics is written, repaired or re-derived: the statement and the proof are the article's own lines, in the article's own notation, and no retained line is substituted or re-flowed.  Retained uncounted as the source-local context the proof needs: lines 265--267.  Nothing was omitted from any retained span, so the package carries no omission record.  No retained line contains a citation, so the wrapper carries no bibliography and no bibliography entry.  No retained line contains a figure environment, an image-inclusion command or drawing code, so no image is delivered and the package declares no asset dependency.  Editorial formatting only, in addition: an AMS \texttt{amsart} preamble that restates the article's own declarations and macros used by the retained lines, namely the environments \texttt{theorem} on separate counters, together with the standard packages those lines need.  The retained environments print in this document as Theorem 1 in order of appearance, whereas the article prints the same items with the numbering of its own full document.  The complete native source of the article as distributed in the arXiv e-print for this exact version is bundled unchanged as \texttt{sources/170381/source0.tex}.
\begin{theorem}\label{zrtversusaht}
$\zrt^2\weiseq\aht$.
\end{theorem}

\begin{theorem}\label{higher-zrtversusaht}
For each $d\in\mathbb N\setminus\{0\}$, we have $\zrt^{d+1}\weiseq\aht^d$.
\end{theorem}

\begin{proof}
The Turing functionals $\Phi_1,\Psi_1$ witnessing that $\zrt^{d+1}\weisleq\aht^d$ are: for a $\mathbb Z$-invariant colouring $c:[\mathbb N]^{d+1}\longrightarrow k$, we let $\Phi_1^c:\mathbb N^d\longrightarrow k$ be given by
\begin{equation*}
\Phi_1^c(y_1,\ldots,y_d)=c(0,y_1,y_1+y_2,\ldots,y_1+y_2+\cdots+y_d).
\end{equation*}
On the other hand, given an infinite set $Y=\{y_n\big|n\in\mathbb N\}$, enumerated increasingly, let $\Psi_1^Y=\{y_1+\cdots+y_n\big|n\in\mathbb N\}$. %Note that, since $c$ is $\mathbb Z$-invariant, we have $\Phi_1^c(y_1,\ldots,y_d)=c(x,x+y_1,x+y_1+y_2,\ldots,x+y_1+\cdots+y_d)$ for every $x\in\mathbb N$.
Now, if $x_0,\ldots,x_d\in\Psi_1^Y$ are distinct elements, with $x_0<\cdots<x_d$, then we have 
$x_i=y_1+\cdots+y_{k_i}$ for each $i$, with
%$y_1+\cdots+y_{k_1},y_1+\cdots+y_{k_2},\ldots,y_1+\cdots+y_{k_d}\in\psi_1^Y$ are $d$ distinct elements (and enumerated increasingly, so that 
$k_0<k_2<\cdots<k_d$. Then, by the $\mathbb Z$-invariance of $c$ we have
\begin{eqnarray*}
c(x_0,\ldots,x_d) & = & c(0,x_1-x_0,x_2-x_0,\ldots,x_d-x_0) \\
 & = & c(0,z_1,z_1+z_2,\ldots,z_1+z_2+\cdots+z_d) \\
 & = & \Phi_1^c(z_1,z_2,\ldots,z_d),
\end{eqnarray*}
where we have defined $z_i=x_i-x_{i-1}=y_{k_{i-1}+1}+\cdots+y_{k_i}$. So the values of $c$ on $(d+1)$-subsets from $\Psi_1^Y$ are exactly the values of $\Phi_1^c$ on $d$-tuples from $Y$ that have the form
\begin{equation*}
(y_{k_0+1}+\cdots+y_{k_1},y_{k_1+1}+\cdots+y_{k_2},\cdots,y_{k_{d-1}+1}+\cdots+y_{k_d}),
\end{equation*}
which are precisely the elements of $\afs^d(Y)$. Hence $\Psi_1^Y$ is $c$-monochromatic if and only if $Y$ is $\Phi_1^c$-monochromatic.

Now, to prove $\aht^d\weisleq\zrt^{d+1}$, the witnessing Turing functionals will be $\Phi_2,\Psi_2$ defined as follows: for $d:\mathbb N^d\longrightarrow k$ we let $\Phi_2^d:[\mathbb N]^{d+1}\longrightarrow k$ be given by $\Phi_2^d(x_0,\ldots,x_d)=d(x_1-x_0,x_2-x_1,\ldots,x_d-x_{d-1})$ (if $x_1<\ldots<x_{d+1}$); it is easy to see that $\Phi_2^d$ is a $\mathbb Z$-invariant colouring. For the definition of $\Psi_2$ we proceed as in the proof of theorem~\ref{zrtversusaht}: given an infinite set $X$, first define $X'=\{x_n\big|n\in\mathbb N\}$ such that $x_0,x_1$ are the two smallest elements of $X$, and $x_{n+1}$ is the least element of $X$ so that $x_{n+1}-x_n>x_n-x_{n-1}$ for each $n\geq 1$. We then let $\Psi_2^X=\{x_{n+1}-x_n\big|n\in\mathbb N\}$. Note that, if $y_j=x_{j+1}-x_j$, then an element $\vec{y}\in\afs^d(\Psi_2^X)$ is of the form
\begin{eqnarray*}
\vec{y} & = & (y_{i_0}+\cdots+y_{i_{1}-1},y_{i_1}+\cdots+y_{i_2-1},\ldots,y_{i_{d-1}}+\cdots+y_{i_d-1}) \\
& = & (x_{i_1}-x_{i_0},x_{i_2}-x_{i_1},\ldots,x_{i_d}-x_{i_{d-1}}), 
\end{eqnarray*}
so that $d(\vec{y})=\Phi_2^d(x_0,\ldots,x_d)$ and hence, if $[X]^{n+1}$ is $\Phi_2^d$-monochromatic this implies that $\afs(\Psi_2^X)$ is $d$-monochromatic.
\end{proof}

\end{document}
